\documentclass[11pt,a4paper]{article}
\usepackage[margin=0.7in, top=0.6in, bottom=0.6in]{geometry}
\usepackage[utf8]{inputenc}
\usepackage[T1]{fontenc}
\usepackage{charter}
\usepackage{fontawesome5}
\usepackage{xcolor}
\usepackage{enumitem}
\usepackage{hyperref}
\usepackage{microtype}

\definecolor{accent}{HTML}{1E3A5F}     % Professional Dark Navy
\definecolor{accentblue}{HTML}{2B6CB0} % Slate Blue
\definecolor{lightgray}{HTML}{F3F4F6}
\definecolor{darktext}{HTML}{1F2937}
\definecolor{bulletline}{HTML}{9CA3AF}
\definecolor{subtletext}{HTML}{4B5563}

\hypersetup{colorlinks=true, linkcolor=accentblue, urlcolor=accentblue}

\setlength{\parindent}{0pt}
\setlength{\parskip}{8pt}

\begin{document}
\thispagestyle{empty}
\color{darktext}

% Encabezado idéntico al CV
\begin{center}
    {\Huge \bfseries\color{accent} SANTIAGO ALEJANDRO OCHOA QUESADA} \\[6pt]
    {\small\color{darktext}
    \faEnvelope\ \href{mailto:santy.alejo4@gmail.com}{santy.alejo4@gmail.com}
    \quad\textcolor{bulletline}{\textbar}\quad
    \faPhone\ +57 320 351 7153
    \quad\textcolor{bulletline}{\textbar}\quad
    \faMapMarker*\ Bogotá, Colombia \\[4pt]
    \faLinkedin\ \href{https://linkedin.com/in/santiago-alejandro-ochoa-quesada-a07a18288/}{LinkedIn}
    \quad\textcolor{bulletline}{\textbar}\quad
    \faGithub\ \href{https://github.com/Sochoaq}{GitHub}
    }
\end{center}
\vspace{2pt}
{\color{bulletline}\hrule height 0.8pt}
\vspace{10pt}

% Metadatos de la carta
\noindent
\begin{minipage}[t]{0.62\textwidth}
    \textbf{\color{accent} Para:}\\
    Equipo de Selección \& Talento TI\\
    \textbf{Grupo Bancolombia}\\
    Colombia
\end{minipage}
\hfill
\begin{minipage}[t]{0.35\textwidth}\raggedleft
    \textbf{\color{accent} Fecha:}\\
    26 de octubre de 2026
\end{minipage}

\vspace{6pt}
\noindent
{\large\bfseries\color{accent} Asunto: Postulación al cargo de Ingeniero/a de Software Nivel 1 (Junior) - TI}
\vspace{2pt}

\noindent Estimado Equipo de Selección de Bancolombia,

Me dirijo a ustedes con gran entusiasmo para presentar mi postulación a las vacantes de \textbf{Ingeniero/a de Software Nivel 1 (Junior)} en el Grupo Bancolombia, particularmente con foco en \textbf{Integración de Datos y Archivos TI} y \textbf{Servicios Transversales}. Como Ingeniero Electrónico de la \textbf{Universidad Nacional de Colombia} (promedio 4.6/5.0) y actual \textbf{Ingeniero de Proyectos} en AM4, combino un riguroso pensamiento analítico con experiencia práctica en desarrollo backend en Python, diseño de bases de datos SQL y construcción de arquitecturas de software confiables.

A lo largo de mi trayectoria técnica, he desarrollado competencias directamente alineadas con los desafíos tecnológicos y estándares de TI del Banco:

\begin{itemize}[label={\color{accentblue}\raisebox{0.2ex}{\tiny\faCircle}}, leftmargin=1.4em, itemsep=3pt, parsep=0pt, topsep=2pt]
    \item \textbf{Integración de Datos y Pipelines (ETL):} Diseño e implementación de flujos automatizados en Python y SQL para la captura, procesamiento, depuración y transformación estructurada de datos masivos y archivos operativos multi-fuente.
    \item \textbf{Desarrollo Backend y APIs:} Creación de microservicios y endpoints REST con FastAPI, modelado entidad-relación y gestión de persistencia relacional con altos estándares de consistencia y seguridad.
    \item \textbf{Calidad, Resiliencia y Despliegue:} Pruebas automatizadas (QA con Pytest), contenedorización en Docker, integración continua con GitHub Actions y monitoreo de variables críticas, logrando un 98\% de confiabilidad en entornos operativos para clientes de alta exigencia.
\end{itemize}

Cuento con sólida formación en redes y seguridad (\textbf{Cisco CCNA y Ethical Hacking}), dominio de \textbf{inglés intermedio-avanzado (B2)} y adaptabilidad a marcos de trabajo colaborativos y ágiles (Scrum). Me motiva profundamente el propósito de Bancolombia de impulsar el desarrollo sostenible y la inclusión financiera a través de tecnología robusta y de vanguardia.

Agradezco de antemano el tiempo dedicado a revisar mi perfil y quedo a su entera disposición para ampliar cualquier información en una entrevista técnica.

\vspace{10pt}
\noindent Cordialmente,

\vspace{12pt}
\noindent
\textbf{\large\color{accent} Santiago Alejandro Ochoa Quesada}\\
{\small\color{subtletext} Ingeniero Electrónico | AI \& Software Engineering}\\
{\footnotesize\color{subtletext} santy.alejo4@gmail.com \;\textbar\; +57 320 351 7153 \;\textbar\; Bogotá, Colombia}

\end{document}
